% !TEX root = ../main.tex
%--- Slide 4: il modello di rumore (circa 1 min 15 s) ------------------------------------
\begin{frame}{Il modello di rumore}
  \begin{itemize}
    \item \textbf{White noise}: spettro piatto, errori indipendenti da un TOA all'altro (fissato)
    \item \textbf{Rumore rosso}: fluttuazioni correlate nel tempo, con più potenza
          alle basse frequenze: derive lente dei residui
    \item Densità spettrale a legge di potenza:
          $S(f) \propto A^2 \left(f/\fyr\right)^{-\gamma}$\\
          $A$: ampiezza alla frequenza $\fyr = 1/\text{anno}$ \quad
          $\gamma$: quanto la potenza è concentrata a bassa frequenza
  \end{itemize}
  \vspace{2mm}
  \centering
  \renewcommand{\arraystretch}{1.15}
  \begin{tabular}{lll}
    \toprule
    Processo & Origine & Dipendenza da $\nu$ radio \\
    \midrule
    Red noise intrinseco (RN) & rotazione della pulsar & acromatico \\
    Variazioni di DM          & densità di elettroni & $\propto \nu^{-2}$ \\
    Variazioni di scattering (SV) & mezzo interstellare & $\propto \nu^{-\alpha}$, \ $\alpha = 4$ \\
    \bottomrule
  \end{tabular}
  \vspace{2mm}
  \begin{itemize}
    \item Modello di riferimento: \textbf{6 parametri liberi}
    \item Il GWB, in una singola pulsar, è un altro rumore rosso acromatico
  \end{itemize}
\end{frame}
